\section{Introduction}
Present baryon abundance retains information about the excess of matter over antimatter, but reconstructing the quantities that coexisted during annihilation requires establishing what additional information a thermal model supplies. Two inventories with the same difference between baryons and antibaryons can contain very different pair populations. The retrospective question must therefore distinguish the observed residual from the total population inferred at a specified temperature.

The literature provides concrete reference points for this comparison. \citet{Satarov2013} describe nearly equal baryon and antibaryon abundances above approximately 50 MeV and a rapid decrease in the antibaryon fraction as the plasma cools below that scale. \citet{Yang2022} investigate the 50--30 MeV window and report 38.2 MeV for the criterion $n_{\bar B}/(n_B-n_{\bar B})=1$. These reference points describe abundance evolution; 50 MeV is not a measurement of the first annihilation event, and the second criterion does not require antibaryons to have disappeared.

The hadron resonance gas (HRG) permits equilibrium populations to be calculated from the rest masses, spins, and charges of particle species. Thermal-FIST lists provide a documented inventory for this calculation \citep{Vovchenko2019,TFLists2026}. In this work, these tabulated properties constitute the signatures used to identify each particle species. No original spectral signals from the early epoch were analyzed. Present cosmological abundance was incorporated as an external constraint on net baryon number.

The objective was to test whether these properties, combined with present abundance and a temperature grounded in the literature, permit separate inference of early baryon and antibaryon populations, quantification of their relationship to the residual, and comparison with published results. Whether the present determines a unique historical reconstruction was also examined. The analysis focused on 50 MeV and the location of the published abundance crossing, retaining the HRG inventory masses and stating the conditions that enable the inference.
